% ECONOMICS_PROBLEM: {"id": "E43-524", "title": "Real-time measurement of the natural interest rate", "jel_code": "E43", "source_id": "OP-0013"}
% !TeX program = xelatex
\documentclass[11pt,a4paper]{article}
\usepackage[margin=23mm]{geometry}
\usepackage{fontspec}
\setmainfont{Latin Modern Roman}
\usepackage{amsmath,amssymb,mathtools}
\usepackage{xcolor,enumitem,booktabs,longtable,array,xurl,hyperref}
\definecolor{muted}{HTML}{53636C}
\hypersetup{colorlinks=true,urlcolor=blue,linkcolor=blue}
\setlength{\parindent}{0pt}
\setlength{\parskip}{5pt}
\newcommand{\R}{\mathbb R}
\newcommand{\E}{\mathbb E}
\newcommand{\N}{\mathbb N}
\newcommand{\ind}{\mathbf 1}
\newcommand{\Var}{\operatorname{Var}}
\newcommand{\Cov}{\operatorname{Cov}}
\newcommand{\supp}{\operatorname{supp}}
\DeclareMathOperator*{\argmin}{arg\,min}
\DeclareMathOperator*{\argmax}{arg\,max}
\newcommand{\entrymeta}[3]{{\sffamily\small\color{muted}\textbf{\texttt{#1}}\quad|\quad #2\par #3\par}}
\newcommand{\Problemref}[1]{\texttt{#1}}
\newcommand{\JELref}[2]{\texttt{#2} (JEL \texttt{#1})}
\begin{document}
\section*{Real-time measurement of the natural interest rate}
\textbf{Problem ID:} \texttt{E43-524}\quad \textbf{JEL:} \texttt{E43}\par
% BEGIN ECONOMICS STATEMENT
\label{problem:OP-0013}
\label{atlas:prob:A3}

\noindent\textbf{Original atlas ID:} A3. \textbf{Formulation:} editorial specification of a research family.

\subsection*{Definitions and assumptions}

Restore the missing symbol in the atlas question as $r_t^*$, the model-dependent real short rate consistent with a specified natural-output benchmark and stable inflation after excluding the declared transitory disturbances. Let $\mathcal D_t^{(v)}$ denote the data vintage actually available at decision date $t$; revised future data are not admissible inputs. A candidate state-space model $M\in\mathcal M$ specifies latent state $s_t$, transition law $s_t\sim K_M(\cdot\mid s_{t-1})$, and observations $d_t\sim H_M(\cdot\mid s_t)$, with $r_t^*=g_M(s_t)$. Parameters, initial-state laws, expectation measures, and revision processes are all part of $M$. Let $\mathcal M_I$ be the model set consistent with the declared identifying restrictions and available-data law.

\subsection*{Formal research question}

Characterize real-time information about $r_t^*$ across $\mathcal M_I$, and evaluate whether it is sufficient for a specific policy decision and loss. Seek uncertainty sets $C_t(\mathcal D_t^{(v)})$ satisfying a declared repeated-sampling coverage criterion, for example
\[
\inf_{M\in\mathcal M_I}\Pr_M\{r_t^*\in C_t(\mathcal D_t^{(v)})\}\ge1-\alpha,
\]
where $\alpha\in(0,1)$ and the model class must be restricted enough for the requirement to be informative. For an admissible action $a$ and bounded loss $L_M(a,r_t^*)$, compare policies by their real-time expected loss and worst-model regret. Quantify separately vintage revisions, filtering uncertainty, parameter uncertainty, and disagreement about what the natural-rate target means.

\subsection*{Interpretation and scope}

A posterior mean from one model and a frequentist confidence region for an identified parameter are different objects. The latent rate itself can retain nonvanishing filtering uncertainty even when model parameters are estimated accurately; an arbitrary requirement of shrinking pointwise intervals would therefore be incorrect. A trend natural rate and a short-run flexible-price rate also need not coincide. ``Sufficiently well'' must be expressed using a policy horizon, feasible actions, and an explicit loss or regret tolerance. Historical smoothing with revised data may explain a past path but cannot demonstrate real-time usefulness. An ensemble should preserve each model's target definition and should report whether disagreements concern measurement or economically different counterfactual rates.

\subsection*{Source audit and status}

Laubach--Williams (2003), Holston--Laubach--Williams (2017), and Del Negro--Giannone--Giannoni--Tambalotti (2017) are verified. The official natural-rate bibliography confirms the 2003 journal article. The Holston et al. publisher record gives pages S59--S75; secondary or official summaries can contain inconsistent pagination. Original link [17] points instead to \emph{Measuring the Natural Rate of Interest Redux}, a distinct 2016 FEDS paper, now listed as [S4]. Its uncertainty discussion supports the motivation but must not be attributed to the 2003 title. The coverage and decision-theoretic targets above are editorial research specifications; no single universally observable $r^*$ exists independently of a model.

\subsection*{References}

\begin{enumerate}

\item[S1] Thomas Laubach and John C. Williams (2003). \emph{Measuring the Natural Rate of Interest}. Review of Economics and Statistics 85(4), 1063--1070. \url{https://www.newyorkfed.org/research/policy/rstar}. \textit{Locator:} official bibliography and description of the Laubach--Williams model. \textit{Establishes:} A latent-state model for natural-rate estimation with real-time uncertainty. \textit{Does not establish:} precise observation or unrestricted identification of the natural rate.

\item[S2] Kathryn Holston, Thomas Laubach, and John C. Williams (2017). \emph{Measuring the Natural Rate of Interest: International Trends and Determinants}. Journal of International Economics 108, supplement 1, S59--S75. \url{https://www.sciencedirect.com/science/article/pii/S0022199617300065}. \textit{Locator:} abstract; cross-economy model. \textit{Establishes:} An international extension and evidence about common natural-rate movements. \textit{Does not establish:} a model-free universal natural interest rate.

\item[S3] Marco Del Negro, Domenico Giannone, Marc P. Giannoni, and Andrea Tambalotti (2017). \emph{Safety, Liquidity, and the Natural Rate of Interest}. Brookings Papers on Economic Activity, Spring, 235--316; also Federal Reserve Bank of New York Staff Report 812. \url{https://www.brookings.edu/articles/safety-liquidity-and-the-natural-rate-of-interest/}. \textit{Locator:} abstract; yield-trend and DSGE analyses. \textit{Establishes:} Evidence linking safe/liquid asset demand and the natural rate. \textit{Does not establish:} real-time point identification independent of model and asset-premium restrictions.

\item[S4] Thomas Laubach and John C. Williams (2016). \emph{Measuring the Natural Rate of Interest Redux}. Finance and Economics Discussion Series 2016-011, Federal Reserve Board. \url{https://www.federalreserve.gov/econres/feds/measuring-the-natural-rate-of-interest-redux.htm}. \textit{Locator:} abstract. \textit{Establishes:} An assessment of robustness and policy consequences of natural-rate uncertainty. \textit{Does not establish:} the identity of the distinct 2003 article.

\end{enumerate}

\subsection*{Common conventions: Source mathematical conventions}

Unless an entry states otherwise, agent and firm sets are finite. State and action spaces are measurable, strategies and outcome maps are measurable, probability laws are specified, and expectations exist for the targets under discussion. Infinite-horizon discount factors lie in $(0,1)$ and discounted objectives are finite. Norms, tolerances and complexity input sizes must be specified for computational comparisons. Constants or primitives that appear locally are inputs, not empirical facts asserted by this catalogue.

\textbf{Identification.} A statistical model $m\in\mathcal M$ implies an observable law $P_m$ and target $\theta(m)$. At law $P$, its identified set is
\[
\Theta(P;\mathcal M)=\{\theta(m):m\in\mathcal M,\ P_m=P\}.
\]
Point identification holds when this set is a singleton, or equivalently when $P_{m_1}=P_{m_2}$ implies $\theta(m_1)=\theta(m_2)$ for all compatible models. Sharp bounds mean bounds attained, or approached when the boundary is not attained, by compatible models. Writing this set is a definition, not a solution: the substantive task is to establish informative restrictions and derive or estimate the set. An unrestricted-confounding model may give only trivial bounds.

\textbf{Inference.} Identification, estimability and valid uncertainty quantification are separate questions. Fix $\alpha\in(0,1)$. A confidence procedure $C_n$ over a declared law class $\mathcal P$ has asymptotic uniform coverage at level $1-\alpha$ when
\[
\liminf_{n\to\infty}\inf_{P\in\mathcal P}\Pr_P\{\theta(P)\in C_n\}\geq1-\alpha.
\]
For set-identified targets the coverage event must be defined separately, for example coverage of the full identified set. If an entry uses identification-indexed models instead of uniquely indexed laws, the same coverage convention is applied over those models. Pointwise limits do not establish uniform validity, and asymptotic validity does not guarantee finite-sample precision. Sample sizes, dependence structures and weak-identification sequences are part of each application.

\textbf{Counterfactuals.} $Y_i(a)$ is a potential outcome under a specified intervention, including its timing, implementation and versions. It is not $Y_i$ conditional on voluntarily choosing $a$. A full intervention vector permits interference; shorthand scalar treatment notation does not silently impose no interference. Assignment, selection, attrition, anticipation and measurement must be specified. A claim about an instrument's local effect is not automatically a population policy effect.

\textbf{Economic equilibrium.} A counterfactual model includes best responses, constraints, feasibility, market clearing, state transitions and expectations consistent with the stated information. When several equilibria exist, either a declared selector is an input or the target is set-valued. Comparisons across policy regimes must use the stated selection convention. Reduced-form relationships are not presumed invariant after an equilibrium-changing intervention.

\textbf{Optimization and welfare.} Optimization is over the stated feasible set. If attainment is not established, $\sup$ or $\inf$ and $\varepsilon$-optimal choices replace an asserted maximizer or minimizer. For example, with value $V=\sup_{a\in\mathcal A}W(a)<\infty$, an $\varepsilon$-optimal set is $\{a\in\mathcal A:W(a)\geq V-\varepsilon\}$ for $\varepsilon>0$. Benefits, costs, risk and health or environmental outcomes can be aggregated only using declared units and conversion weights. Interpersonal utility weights, the treatment of fiscal transfers and foreign profits, discounting, and the behavioral welfare criterion are explicit inputs. A comparative-static derivative additionally requires existence and appropriate differentiability of the selected counterfactual.

\textbf{Sources.} Local labels [S1], [S2], and so forth restart in every question. Locators identify the relevant abstract, model, section or result at the level actually checked; a bibliographic or abstract-level verification is not represented as inspection of an entire proof. Working papers are distinguished from later publications and changing versions. Source summaries are paraphrased. All URL checks and editorial formulations refer to the audit date stated above.
% END ECONOMICS STATEMENT
\end{document}
